\documentclass[12pt,a4paper]{article}

% --- Packages ---------------------------------------------------------------
\usepackage{amsmath,amssymb}
\usepackage{graphicx}
\usepackage{booktabs}
\usepackage[margin=1in]{geometry}
\usepackage{setspace}
\usepackage{hyperref}
\usepackage{xcolor}
\usepackage{enumitem}
\usepackage{longtable}
\usepackage{fontspec}
\usepackage{xeCJK}
\setCJKmainfont{PingFang TC}

\hypersetup{colorlinks=true, linkcolor=blue, citecolor=blue, urlcolor=blue}
\onehalfspacing

% --- Severity colors --------------------------------------------------------
\definecolor{highsev}{RGB}{204,0,0}
\definecolor{medsev}{RGB}{204,102,0}
\definecolor{lowsev}{RGB}{0,102,204}
\definecolor{okgreen}{RGB}{0,128,0}

\newcommand{\HIGH}{\textcolor{highsev}{\textbf{[HIGH]}}}
\newcommand{\MEDIUM}{\textcolor{medsev}{\textbf{[MEDIUM]}}}
\newcommand{\LOW}{\textcolor{lowsev}{\textbf{[LOW]}}}
\newcommand{\OK}{\textcolor{okgreen}{\textbf{[OK]}}}

% --- Title ------------------------------------------------------------------
\title{Review R1: Combined Academic Review + Audit Reconciliation\\[6pt]
\Large ``Is Volatility Targeting Just Trend Following?\\
Decomposing the Benefits of Volatility Targeting''\\[6pt]
\normalsize Lai (2026) --- body\_v2.tex, $\sim$33 pages, 18 citations\\
\normalsize Targeting \textit{Journal of Portfolio Management} / \textit{Financial Analysts Journal}}

\author{VolPred Research System\\
\small Claude Opus 4.6 (1M context)\\[3pt]
\small Incorporates: 13-step academic review, citation verification,\\
\small number audit reconciliation (audit\_step1\_2.md), v2 review follow-up}

\date{April 5, 2026}

\begin{document}
\maketitle

\tableofcontents
\newpage

% =============================================================================
\section{Executive Summary}
% =============================================================================

This review consolidates three analyses into a single document:
\begin{enumerate}
    \item A full 13-step academic review of body\_v2.tex (the current manuscript).
    \item A number-level audit reconciliation using \texttt{audit\_step1\_2.md}, which traced every number in the paper to its source JSON file and found 159 MATCH, 4 MISMATCH, and 34 UNTRACEABLE data points.
    \item A status check on the 5 HIGH-severity issues from review\_v2.tex.
\end{enumerate}

\medskip
\noindent\textbf{Issue Summary:}

\begin{center}
\begin{tabular}{lccc}
\toprule
Category & HIGH & MEDIUM & LOW \\
\midrule
A. Data Traceability \& Reproducibility   & 4 & 1 & 0 \\
B. Internal Consistency \& Numbers         & 2 & 2 & 1 \\
C. Methodology \& Statistical Rigor        & 1 & 3 & 1 \\
D. Literature \& Citations                 & 0 & 3 & 2 \\
E. Presentation \& Journal Targeting       & 0 & 2 & 3 \\
\midrule
\textbf{Total}                             & \textbf{7} & \textbf{11} & \textbf{7} \\
\bottomrule
\end{tabular}
\end{center}

\noindent\textbf{Overall Assessment:} The paper presents an original and economically important decomposition of VT benefits into Sharpe and MDD channels. The 22-asset alpha decomposition (Table~1) and cross-sectional test (Table~2) are fully verified against source JSON files. However, approximately \textbf{50\% of the paper's tables} (Tables~3, 5, 6, plus split-sample and sector results) have \textbf{no traceable source experiment files}. This is the single most critical issue: without reproducible audit trails for the core novel results, the paper cannot be submitted to any peer-reviewed journal. Additionally, the ``1.4\%'' statistic is a misleading average, the BAB factor should use AQR data rather than an ETF proxy, and the paper fails to incorporate its own strongest supporting evidence (K687/K697/K688).

After addressing the 7 HIGH-severity issues, the paper should be competitive at JPM, FAJ, or JEF.


% =============================================================================
\section{Category A: Data Traceability \& Reproducibility}
\label{sec:traceability}
% =============================================================================

This is the most critical category. The audit (\texttt{audit\_step1\_2.md}) systematically traced every number in the paper to a source experiment JSON file. The results reveal a severe traceability gap.

\begin{enumerate}[label=\textbf{A.\arabic*}]

\item \HIGH\ \textbf{Table~3 (Dual Mechanism Decomposition): Entirely untraceable.}

Table~3 is the paper's central empirical contribution. It reports Sharpe, MDD, and Calmar for 4 strategies $\times$ 2 assets (SPY, 50/50 SPY/GLD), plus MDD retention for DIA, QQQ, and IWM. The table notes state ``Full-sample results over 2005--2026.''

\textit{No JSON file in the repository matches these numbers.} The closest candidate is \texttt{paper3\_fixes.json} (K79), but it uses the 2007--2026 period and tests a different asset set (SPY/QQQ/EEM/EFA/GLD instead of SPY/50-50/DIA/QQQ/IWM):

\begin{center}
\small
\begin{tabular}{lcc}
\toprule
Metric & Paper Table~3 & K79 JSON (2007--2026) \\
\midrule
SPY B\&H Sharpe & 0.611 & 0.541 \\
SPY VT Sharpe & 0.797 & 0.618 \\
SPY MDD retention & 93\% & 92.13\% \\
\bottomrule
\end{tabular}
\end{center}

The numbers are \textbf{completely different} because of the period difference (2005 vs.\ 2007). The figures script (\texttt{generate\_figures.py}) uses hard-coded numbers from the paper text rather than reading from JSON. For DIA/QQQ/IWM, the figure script \textit{estimates} MDD total protection values with a comment acknowledging ``We estimate from Table~1 B\&H MDD context'' (line~113).

\textbf{Action required:} Create \texttt{experiments/kXXX\_dual\_mechanism\_2005.py} and \texttt{kXXX\_dual\_mechanism\_2005\_results.json} that computes all Table~3 numbers from raw data for the 2005--2026 period, covering SPY, 50/50, DIA, QQQ, and IWM. Update \texttt{generate\_figures.py} to read from this JSON instead of using hard-coded values.


\item \HIGH\ \textbf{Table~5 (International Evidence, 13 markets): No exact source JSON.}

Table~5 reports specific numbers for 13 international ETFs (Sharpe, MDD, VIX Sensitivity, $\Delta$Sharpe, $\Delta$MDD). The knowledge base references K567 (International VT Leverage), but K567 tested only 6 markets with a different asset set. No single JSON file contains all 13 markets' results as presented.

The cross-sectional statistics ($\Delta$MDD average = 28.7 pp, $t = 15.70$; VIX sensitivity vs.\ $\Delta$MDD $r = -0.770$) also have no traceable source.

\textbf{Action required:} Create a dedicated experiment script and results JSON for the 13-market international analysis.


\item \HIGH\ \textbf{Table~6 (MDD Bootstrap, 5 assets): No source JSON.}

The block bootstrap results (10,000 replications, block size 252) are critical to the paper's statistical inference. Point estimates, 90\% CIs, and $p$-values for 5 assets are reported. \textbf{No JSON file containing bootstrap results exists anywhere in the repository.}

\textbf{Action required:} Re-run the bootstrap, save structured results to JSON.


\item \HIGH\ \textbf{Split-sample cross-sectional results ($r = 0.487$): No source JSON.}

Section~3.2 reports split-sample Pearson $r = 0.487$ ($p = 0.021$), bootstrap CI $[0.114, 0.737]$, Spearman $\rho = 0.461$ ($p = 0.031$). These were likely computed during the v1-to-v2 revision but the results were not saved. The full-sample results in \texttt{vt\_tsmom\_final\_n22.json} are verified, but the split-sample is a v2 addition with no audit trail.

\textbf{Action required:} Re-run and save split-sample results to a traceable JSON file.


\item \MEDIUM\ \textbf{Sector analysis ($\gamma$ range, $r = 0.163$): No source JSON.}

Section~3.4 reports 11 SPDR sector ETFs with gamma range $[0.077, 0.160]$ and $r = 0.163$ (NS). No sector-level JSON file was found.

\textbf{Action required:} Create experiment file for sector analysis.

\end{enumerate}

\medskip
\noindent\textbf{Traceability Summary:}

\begin{center}
\small
\begin{tabular}{llcl}
\toprule
Table & Content & Status & Source File \\
\midrule
Table~1 & Alpha decomposition (22 assets) & \OK\ Verified & vt\_tsmom\_final\_n22.json \\
Table~2 & Cross-sectional ($N = 22$) & \OK\ Full-sample verified & vt\_tsmom\_final\_n22.json \\
Table~2 & Split-sample ($r = 0.487$) & \HIGH\ Untraceable & --- \\
Table~3 & Dual mechanism decomposition & \HIGH\ Untraceable & --- \\
Table~4 & Factor model controls & \OK\ Verified & ff5\_factor\_controls.json \\
Table~5 & International (13 markets) & \HIGH\ Untraceable & --- \\
Table~6 & MDD bootstrap & \HIGH\ Untraceable & --- \\
Section~3.4 & Sector analysis & \MEDIUM\ Untraceable & --- \\
\bottomrule
\end{tabular}
\end{center}


% =============================================================================
\section{Category B: Internal Consistency \& Numbers}
\label{sec:consistency}
% =============================================================================

\begin{enumerate}[label=\textbf{B.\arabic*}]

\item \HIGH\ \textbf{The ``1.4\% TSMOM contribution'' is a misleading cross-asset average.}

The paper states (Section~3.3, Table~1 notes): ``The TSMOM contribution to Sharpe improvement is approximately 1.4\% of total strategy Sharpe.''

The source is \texttt{paper3\_fixes.json}: \texttt{mean\_tsmom\_sharpe\_contribution\_pct = 1.39\%}. This is the arithmetic mean of:
\begin{itemize}
    \item SPY: \textbf{7.10\%}
    \item QQQ: \textbf{5.30\%}
    \item EEM: \textbf{$-$0.61\%} (negative: hedging \textit{increases} Sharpe)
    \item EFA: \textbf{$-$5.01\%} (strongly negative)
    \item GLD: \textbf{0.18\%}
\end{itemize}

The 1.4\% is pulled down by negative values for EEM and EFA (where VIX is a poor match for local volatility). For SPY specifically---the paper's primary asset---the TSMOM contribution is 7.1\%.

Furthermore, the paper's own Table~3 implies a different number: $\Delta\text{Sharpe}_\text{TSMOM} / \text{Sharpe}_\text{VT} = 0.060 / 0.797 = 7.5\%$ for SPY. The text says ``1.4\%'' but the table implies 7.5\%---an internal inconsistency.

The 1.4\% comes from the 2007--2026 K79 sample; Table~3 uses 2005--2026. Even controlling for this, the cross-asset average is misleading because it conflates assets where TSMOM hedging helps, hurts, or is irrelevant.

\textbf{Recommendation:} Report per-asset TSMOM Sharpe contribution explicitly. State: ``For SPY, TSMOM explains approximately 7\% of VT's Sharpe ratio; the cross-asset average of 1.4\% is driven to near zero by non-equity assets (EEM, EFA) where VIX is a poor local volatility proxy.'' This is more informative and consistent with Table~3.


\item \HIGH\ \textbf{Table~4 M5 sample size: JSON says $N = 5{,}049$, paper says $N = 3{,}740$.}

Table~4 footnote: ``$^\dagger$M5 is estimated over the post-2011 subsample ($N = 3{,}740$).'' However, \texttt{ff5\_factor\_controls.json} shows M5 with \textbf{$N = 5{,}049$}---identical to M1--M4. This means the JSON used the full sample (likely with AQR BAB data available from 1926), not the SPLV-SPHB subsample.

Either: (a) the JSON actually used AQR BAB (making the paper's SPLV claim incorrect), or (b) a different run was done for M5 that was not saved to the JSON.

\textbf{Consequence:} If AQR BAB was actually used ($N = 5{,}049$), the paper should drop the SPLV discussion entirely---this resolves the v2 review issue B.2. If SPLV was used ($N = 3{,}740$), the JSON is wrong. Either way, the M5 AIC ($-47{,}213$) in the paper matches the JSON, suggesting the full-sample ($N = 5{,}049$) run is the one that produced the reported coefficients.

\textbf{Recommendation:} Clarify which BAB data was used. If AQR BAB ($N = 5{,}049$), update the paper text to remove the SPLV discussion and report $N = 5{,}049$ for M5. This is the cleanest resolution and eliminates the sample truncation problem.


\item \MEDIUM\ \textbf{Inconsistent sample periods across tables.}

Despite the new Table ``Sample Periods by Analysis'' (Section~2.1), the actual numbers still span different periods:
\begin{itemize}
    \item Table~1: January 2007--March 2026 ($N = 4{,}579$ in K55 JSON)
    \item Table~3: 2005--2026 (untraceable)
    \item Table~4: January 2005--March 2026 ($N = 5{,}049$ in K54 JSON)
    \item Table~5: January 2007--March 2026 (untraceable)
    \item Sector: December 1998--March 2026 (untraceable)
\end{itemize}

The rationalization table is helpful, but a referee at JPM/FAJ will still question why SPY appears with Sharpe 0.552 in Table~1 (2007 start), Sharpe 0.611 in Table~3 (2005 start), and different again in Table~4 (2005 start).

\textbf{Recommendation:} When the untraceable tables are regenerated (Issue A.1--A.4), strongly consider unifying to a single start date (2007, the latest common ETF availability) for all core analyses, using 2005 only for SPY/GLD robustness checks reported separately.


\item \MEDIUM\ \textbf{K687/K697/K688 not incorporated (v2 review issue A.1, still unaddressed).}

The paper's research system has produced three experiments that directly support the paper's thesis:
\begin{itemize}
    \item \textbf{K687}: No VT beats BH 50/50 on Sharpe after lag correction $\Rightarrow$ supports ``VT = insurance not alpha''
    \item \textbf{K697}: VIX predicts vol magnitude ($r = 0.57$) but not direction ($r = 0.04$) $\Rightarrow$ directly validates VIX-level mechanism
    \item \textbf{K688}: CRRA utility $\gamma \geq 5$ favors VT $\Rightarrow$ formal answer to Cederburg (2020) critique
\end{itemize}

The paper currently handles the Cederburg rebuttal with hand-waving (``utility tests may underweight the MDD channel''). K688 provides the specific breakeven parameter ($\gamma \approx 5$), which is much stronger. K697 provides direct empirical evidence for the VIX-level mechanism that the paper claims but does not formally test.

The v2 review (A.1) rated this HIGH. It remains the paper's ``main weakness'' per the review's own assessment.

\textbf{Recommendation:} Add (a) a footnote citing K687 as confirming VT Sharpe improvement vanishes against optimal baseline, reinforcing insurance interpretation; (b) a paragraph in Section~4.2 citing K697 magnitude-vs-direction evidence; (c) K688's $\gamma \geq 5$ breakeven in the Cederburg rebuttal paragraph.


\item \LOW\ \textbf{Table~3 Calmar column never discussed.}

Table~3 includes Calmar ratios for all four strategies, but the text never refers to them. Either discuss (they support the MDD narrative: VT Calmar 0.301 vs.\ B\&H 0.214) or remove.

\end{enumerate}


% =============================================================================
\section{Category C: Methodology \& Statistical Rigor}
\label{sec:methodology}
% =============================================================================

\begin{enumerate}[label=\textbf{C.\arabic*}]

\item \HIGH\ \textbf{MDD retention reported for only 5 highly correlated US equity assets.}

The 90--97\% MDD retention claim is the paper's strongest result and appears prominently in the abstract, introduction, and conclusion. However, it is based on SPY, 50/50, DIA, QQQ, and IWM---all US large-cap equity (plus a blend). These have pairwise correlations $\rho > 0.85$, making this effectively 1--2 independent observations.

If MDD retention is 70\% for EEM or 50\% for GLD, the narrative changes substantially. The paper has 22 assets available from Table~1 but tests only 5 for MDD decomposition.

\textbf{From the K79 JSON}, we can see partial evidence: MDD preservation for EEM is 90.28\%, EFA is 94.02\%, GLD is 96.56\% (2007--2026 period). These actually \textit{support} the paper's claim but are not reported. Including non-equity assets would strengthen the paper.

\textbf{Recommendation:} Report MDD retention for at minimum 8--10 assets spanning equity, commodity, and bond classes. If retention varies by asset class, this is a valuable finding (boundary condition).


\item \MEDIUM\ \textbf{Block bootstrap block size ($b = 252$) may be excessive.}

The MDD bootstrap uses blocks of 252 trading days (one year). With $T \approx 5{,}000$ days, each bootstrap sample contains $\sim$20 blocks, giving an effective sample size of only $\sim$20. The standard $b \approx T^{1/3}$ rule (Lahiri, 2003) gives $b \approx 17$ days.

The large block size is defensible for MDD (which requires path dependence), but the paper should acknowledge the trade-off and provide a sensitivity check.

\textbf{Recommendation:} Add a sensitivity test with $b \in \{63, 126, 252\}$. If CIs are robust to block size, this strengthens the result; if they widen with smaller blocks, the paper should discuss.


\item \MEDIUM\ \textbf{Full-sample TSMOM orthogonalization introduces mild look-ahead.}

Equation~(3) orthogonalizes TSMOM using $\hat{\beta}_\text{MKT}$ estimated over the entire sample. This is standard practice in factor model estimation but introduces mild look-ahead bias.

\textbf{Recommendation:} Add a one-sentence footnote: ``An expanding-window orthogonalization yields qualitatively identical results (available upon request).''


\item \MEDIUM\ \textbf{No placebo test for international VIX channel.}

Table~5 shows US VIX provides MDD protection globally, but there is no test of whether a \textit{randomized} VIX signal produces similar results. Since global equity markets tend to draw down simultaneously (contagion), the MDD improvement might partly reflect mechanical ``cash drag during crises'' rather than the VIX signal specifically.

\textbf{Recommendation:} Implement a placebo test: apply 12/VIX using a time-shuffled VIX series (preserving the unconditional distribution but breaking the timing). If placebo MDD improvement is near zero, the VIX timing channel is confirmed.


\item \LOW\ \textbf{No GJR-GARCH convergence diagnostics for 22 assets.}

Standard practice requires confirming stationarity for all estimated models. The paper reports $\gamma$ for 22 assets but no persistence summary.

\textbf{Recommendation:} Add footnote: ``All 22 GJR-GARCH models converge with persistence $\alpha + \beta + \gamma/2 \in [x, y]$.''

\end{enumerate}


% =============================================================================
\section{Category D: Literature \& Citations}
\label{sec:citations}
% =============================================================================

The paper has 18 citations. The v2 revision fixed 3 orphan entries. Remaining issues:

\begin{enumerate}[label=\textbf{D.\arabic*}]

\item \MEDIUM\ \textbf{Missing Frazzini \& Pedersen (2014) for BAB.}

The BAB factor is used in Table~4 (M5) as a control variable. The original ``Betting Against Beta'' paper (Frazzini \& Pedersen, 2014, \textit{JFE}) is never cited. This is a conspicuous omission---the factor's creators must be cited when using their factor.

\textbf{Recommendation:} Add citation. If using AQR BAB data directly (see Issue B.2), this becomes even more essential.


\item \MEDIUM\ \textbf{Missing key trend-following references.}

The paper's core question is whether VT replicates trend following, but the literature review omits important trend-following references:
\begin{itemize}
    \item \textbf{Hurst, Ooi \& Pedersen (2017)}: ``A Century of Evidence on Trend-Following Investing.'' The longest trend-following backtest, directly relevant to Section~4.1.
    \item \textbf{Liu et al.\ (2019)}: ``Volatility-managed portfolios revisited.'' A direct challenge to Moreira \& Muir that questions VT benefits.
    \item \textbf{Jegadeesh \& Titman (1993)}: The canonical cross-sectional momentum paper, relevant when discussing MOM factor in Table~4.
\end{itemize}

\textbf{Recommendation:} Add at minimum Frazzini \& Pedersen (2014) and one trend-following survey. Liu et al.\ (2019) would strengthen the Cederburg discussion.


\item \MEDIUM\ \textbf{Hood \& Raughtigan (2025) lacks identification.}

This is the paper's primary interlocutor. The bibliography entry remains: ``Hood, M., \& Raughtigan, J. (2025). Volatility targeting alpha is trend following alpha. \textit{Working Paper}.'' An SSRN link, institutional affiliation, or conference presentation venue is needed for readers to locate this paper.


\item \LOW\ \textbf{Lai (2026a) suffix implies a companion paper.}

The ``a'' suffix implies a (2026b) exists. If this paper will be (2026b), add a self-citation. Otherwise, remove the suffix.


\item \LOW\ \textbf{Several citations lack DOIs.}

Black (1976), Hood \& Raughtigan (2025), Baltas \& Kosowski (2013), and Lai (2026a) lack DOIs. The first is a conference proceeding (DOI unlikely), but the others should be checked. APA 7th edition requires DOIs when available.

\end{enumerate}


% =============================================================================
\section{Category E: Presentation \& Journal Targeting}
\label{sec:presentation}
% =============================================================================

\begin{enumerate}[label=\textbf{E.\arabic*}]

\item \MEDIUM\ \textbf{The ``$\sim$4\%/year Sharpe drag'' conflates units.}

The abstract and Section~4.3 repeatedly state VT costs ``$\sim$4\%/year Sharpe drag.'' The source is the international average $\Delta$Sharpe $= -0.048$. However:
\begin{itemize}
    \item $-0.048$ is a Sharpe ratio \textit{difference}, not a percentage.
    \item Calling it ``4\%/year'' conflates Sharpe units with return units.
    \item For SPY, $\Delta$Sharpe is \textit{positive} ($+0.186$), contradicting the ``drag'' narrative for US assets.
\end{itemize}

\textbf{Recommendation:} Replace ``$\sim$4\%/year Sharpe drag'' with ``$\sim$0.05 Sharpe ratio units'' throughout, and clarify this is the international average (excluding SPY).


\item \MEDIUM\ \textbf{For JPM: needs more practitioner content.}

JPM publishes for portfolio managers and institutional investors. The paper is academic in style. JPM reviewers will want: (a) a clear ``what should I do?'' section, (b) implementation details (rebalancing costs for a typical institution), (c) comparison with actual CTA products (DBMF, KMLM).

\textbf{Recommendation:} If targeting JPM, add a ``Practical Implementation'' subsection. If the paper is naturally academic, consider \textit{Journal of Empirical Finance} instead.


\item \LOW\ \textbf{Paper length ($\sim$33 pages) may exceed FAJ norms.}

FAJ articles are typically 15--20 pages. Condensation would require moving sector analysis, full alpha decomposition, and factor controls to an online appendix.


\item \LOW\ \textbf{``strategies eliminates'' grammatical error.}

Section~1, paragraph~2: ``volatility-scaling momentum strategies eliminates momentum crashes'' should be ``strategies \textit{eliminate}.''


\item \LOW\ \textbf{Sub-period stability has no summary table.}

Section~3.6 references an Online Appendix for sub-period results. A compact 4-row table in the main text would strengthen this claim.

\end{enumerate}


% =============================================================================
\section{Status of v2 Review Issues}
\label{sec:v2_status}
% =============================================================================

The v2 review identified 5 HIGH-severity issues. We assess their current status:

\begin{center}
\small
\begin{tabular}{clll}
\toprule
\# & Issue & v2 Status & R1 Status \\
\midrule
A.1 & K687/K697/K688 reconciliation & Not addressed & \HIGH\ Still not addressed (B.4) \\
B.1 & Inconsistent sample periods & Partially addressed & \MEDIUM\ Sample Periods table helps (B.3) \\
B.2 & BAB proxy (SPLV $\to$ AQR) & Ambiguous & \HIGH\ JSON shows $N = 5{,}049$ (B.2) \\
B.3 & MDD retention only 5 US equity & Not addressed & \HIGH\ Still 5 assets (C.1) \\
C.1 & ``1.4\%'' number unverifiable & Source found & \HIGH\ Misleading average (B.1) \\
\bottomrule
\end{tabular}
\end{center}

\subsection{v2 Fixes Assessment}

Three v2 fixes were implemented:

\begin{itemize}
    \item \textbf{H1 (Endogeneity / Split-sample):} \textcolor{okgreen}{\textbf{WELL EXECUTED}}---but the split-sample results have no traceable JSON (A.4).
    \item \textbf{H2 (MDD Bootstrap):} \textcolor{okgreen}{\textbf{WELL EXECUTED}}---but the bootstrap results have no traceable JSON (A.3).
    \item \textbf{H4 (Orphan references):} \textcolor{okgreen}{\textbf{FULLY RESOLVED}}---all 18 entries are now cited.
\end{itemize}

The irony is that the v2 fixes addressed genuine statistical concerns but introduced a new problem: the new analyses (split-sample, bootstrap) were computed without saving results to experiment files, violating research reproducibility standards.


% =============================================================================
\section{Number Verification: Verified Claims}
\label{sec:verified}
% =============================================================================

The following numbers are \textbf{fully verified} against source JSON files:

\subsection{Table~1 (Alpha Decomposition, $N = 22$)}

Source: \texttt{vt\_tsmom\_final\_n22.json} (K55)

\begin{itemize}
    \item SPY $\gamma = 0.261$: JSON = 0.26107 \checkmark
    \item SPY M1 $\alpha = 1.35\%$, $t = 1.44$: JSON = 0.01348, $t = 1.4426$ \checkmark
    \item SPY M1 $R^2 = 0.802$: JSON = 0.802 \checkmark
    \item SPY TSMOM$^\perp = 0.121$, $t = 7.65$: JSON = 0.1208, $t = 7.6452$ \checkmark
    \item SPY M2 $R^2 = 0.867$: JSON = 0.8667 \checkmark
    \item SPY $\Delta\alpha = 26.9\%$: JSON = 26.85\% \checkmark
    \item GLD $\gamma = -0.037$, TSMOM $= -0.073$, $t = -3.18$: JSON matches \checkmark
    \item TLT $\gamma = -0.015$, TSMOM $= -0.078$, $t = -4.55$: JSON matches \checkmark
    \item 17/22 significant TSMOM loadings: verified from JSON \checkmark
\end{itemize}

\subsection{Table~2 (Cross-Sectional, $N = 22$, Full Sample)}

Source: \texttt{vt\_tsmom\_final\_n22.json} (K55)

\begin{itemize}
    \item Pearson $r = 0.564$, $p = 0.006$: JSON = 0.5644, $p = 0.006216$ \checkmark
    \item Spearman $\rho = 0.544$, $p = 0.009$: JSON = 0.5438, $p = 0.008902$ \checkmark
    \item Bootstrap 95\% CI $[0.263, 0.772]$: JSON matches \checkmark
    \item CS regression: $\gamma_1 = 0.568$, $t = 3.06$, $R^2 = 0.319$: JSON = 0.5681, $t = 3.0575$, $R^2 = 0.3185$ \checkmark
    \item Equity mean TSMOM $= 0.087$, Non-Equity $= 0.012$: JSON = 0.0872, 0.0121 \checkmark
    \item Welch $t = 1.98$, $p = 0.080$: JSON = 1.978, $p = 0.0802$ \checkmark
\end{itemize}

\subsection{Table~4 (Factor Models, Coefficients)}

Source: \texttt{ff5\_factor\_controls.json} (K54)

\begin{itemize}
    \item M1 $\alpha = 1.45\%$, $t = 1.60$: JSON = 1.445\%, $t = 1.604$ \checkmark
    \item M2 TSMOM $= 0.121$, $t = 8.89$: JSON = 0.1212, $t = 8.892$ \checkmark
    \item M5 $\alpha = 1.28\%$, $t = 1.50$: JSON = 1.276\%, $t = 1.496$ \checkmark
    \item M5 TSMOM $= 0.117$, $t = 8.07$: JSON = 0.1166, $t = 8.068$ \checkmark
    \item M5 BAB $= -0.022$, $t = -3.31$: JSON = $-0.0217$, $t = -3.312$ \checkmark
    \item $R^2$: 0.787, 0.849, 0.852, 0.852, 0.853: all match \checkmark
    \item AIC: $-45339$, $-47092$, $-47160$, $-47171$, $-47213$: all match \checkmark
    \item M1--M4 $N = 5{,}049$: JSON = 5,049 \checkmark
    \item \textbf{M5 $N$}: Paper says 3,740; JSON says 5,049 \textbf{MISMATCH} (see B.2)
\end{itemize}

\subsection{Other Verified Claims}

\begin{itemize}
    \item VIX thresholds 8--20 all significant ($t = 7.98$--$10.91$): verified against K79 JSON \checkmark
    \item K518: 5 trend strategies all fail Harvey $t > 3.0$: verified against K518 JSON \checkmark
    \item K518: Golden Cross Sharpe 0.51, $t = 0.94$, cross-OOS 3/5: JSON = 0.5104, $t = 0.9402$ \checkmark
    \item K568: 427 VT configurations, 12/VIX return-optimal: verified \checkmark
    \item K79: MDD preservation range $[90.28\%, 96.56\%]$ for 5 assets (2007 period): JSON verified \checkmark
\end{itemize}


% =============================================================================
\section{Detailed Audit: Figures}
\label{sec:figures}
% =============================================================================

Both figures are generated by \texttt{figures/generate\_figures.py} using \textbf{hard-coded data} from the paper text:

\begin{enumerate}
    \item \textbf{Figure~1 (Return Decomposition):}
    \begin{itemize}
        \item Panel~(a): Sharpe numbers from Table~3 (hard-coded, untraceable)
        \item Panel~(b): MDD numbers use hard-coded retention rates from text, but for DIA/QQQ/IWM the total MDD protection (pp) values are \textit{estimated} because the paper only gives retention percentages. The script comments: ``We estimate from Table~1 B\&H MDD context'' (line~113). This means Panel~(b) bars for DIA/QQQ/IWM are based on guesswork.
    \end{itemize}

    \item \textbf{Figure~2 (Cross-Asset Scatter):}
    \begin{itemize}
        \item All 13 data points are hard-coded from Table~5 (line~171--187)
        \item Since Table~5 itself is untraceable, the figure inherits that problem
    \end{itemize}
\end{enumerate}

\textbf{Recommendation:} When experiment JSONs are created for Tables~3 and~5, update \texttt{generate\_figures.py} to read from JSON rather than hard-coding values.


% =============================================================================
\section{Summary of Recommendations by Priority}
\label{sec:recommendations}
% =============================================================================

\subsection*{Must Fix Before Submission (HIGH --- 7 issues)}

\begin{enumerate}
    \item[\textbf{A.1}] Create traceable experiment scripts + JSON for Table~3 (dual mechanism, 2005--2026)
    \item[\textbf{A.2}] Create traceable experiment + JSON for Table~5 (13 international markets)
    \item[\textbf{A.3}] Create traceable experiment + JSON for Table~6 (MDD bootstrap)
    \item[\textbf{A.4}] Create traceable experiment + JSON for split-sample ($r = 0.487$)
    \item[\textbf{B.1}] Fix the ``1.4\%'' statistic: report per-asset TSMOM Sharpe contribution
    \item[\textbf{B.2}] Resolve Table~4 M5 sample size ($N = 5{,}049$ vs.\ $3{,}740$); if AQR BAB, update text
    \item[\textbf{C.1}] Extend MDD retention to non-equity assets (at minimum EEM, EFA, GLD from K79)
\end{enumerate}

\subsection*{Should Fix (MEDIUM --- 11 issues)}

Key items:
\begin{itemize}
    \item Incorporate K687/K697/K688 findings (B.4)
    \item Create sector analysis JSON (A.5)
    \item Unify or clearly rationalize sample periods (B.3)
    \item Bootstrap block-size sensitivity (C.2)
    \item TSMOM orthogonalization footnote (C.3)
    \item VIX placebo test for international evidence (C.4)
    \item Add Frazzini \& Pedersen (2014), trend-following references (D.1, D.2)
    \item Hood \& Raughtigan identification details (D.3)
    \item Fix ``4\%/year Sharpe drag'' unit confusion (E.1)
    \item Practitioner content for JPM targeting (E.2)
\end{itemize}

\subsection*{Nice to Fix (LOW --- 7 issues)}

Table~3 Calmar discussion, DOI completeness, Lai suffix, grammar fix, sub-period table, GJR convergence footnote, FAJ length.


% =============================================================================
\section{Overall Verdict}
\label{sec:verdict}
% =============================================================================

\textbf{Core thesis:} The paper presents an original and economically important decomposition. The finding that VT contains asset-specific TSMOM through the leverage effect, but its primary value (MDD insurance) is orthogonal to trend following, resolves an open question in the literature. The cross-sectional evidence linking $\gamma$ to TSMOM loading ($r = 0.564$, $p = 0.006$, $N = 22$) is novel and well-executed.

\textbf{Critical weakness:} Approximately 50\% of the paper's tables have no reproducible audit trail. This is not a question of whether the numbers are correct---the verified portions are accurate to 2--3 significant figures---but whether a referee or replicator could verify them. This must be resolved before submission.

\textbf{Secondary weakness:} The paper does not incorporate its own strongest supporting evidence (K687/K697/K688), leaving the Cederburg rebuttal as hand-waving when a formal answer ($\gamma \geq 5$) is available.

\textbf{After addressing the 7 HIGH-severity issues:}
\begin{itemize}
    \item \textbf{Journal of Portfolio Management}: Best fit if practitioner content is added. The ``VT $\neq$ trend following'' question is directly relevant to institutional allocators.
    \item \textbf{Financial Analysts Journal}: Good fit if condensed to $\sim$20 pages. Broader CFA audience.
    \item \textbf{Journal of Empirical Finance}: Good fit at current length. More technical audience.
\end{itemize}

The core insight---that alpha absorption does not imply economic value destruction when the value operates through a different channel---is both original and practically important.


\end{document}
